\documentclass[10pt,a4paper]{amsart}
\usepackage[margin=19mm]{geometry}
\usepackage{amsmath,amssymb,mathtools}
\usepackage[scr=rsfs,cal=euler]{mathalfa}
\usepackage{xfrac}
\usepackage[unicode,hidelinks,bookmarks=false]{hyperref}
\newcommand{\ie}{i.e.\ }
\newcommand{\cf}{cf.\ }
\newcommand{\resp}{respectively\ }
\newcommand{\Subsection}{\subsection}
\providecommand{\nobreakdash}{\nobreak\hbox{-}}
\setlength{\emergencystretch}{2em}
\multlinegap0pt
\usepackage{bm}
\usepackage{bbm}
\usepackage{dsfont}
\usepackage{xspace}
\usepackage{cancel}
\usepackage{mathtools}
\usepackage{amsthm}
\makeatletter
\@ifundefined{theorem}{\newtheorem{theorem}{Theorem}}{}
\@ifundefined{claim}{\newtheorem{claim}{Claim}}{}
\@ifundefined{lemma}{\newtheorem{lemma}[theorem]{Lemma}}{}
\makeatother
\title[Borel graphs generated by commuting functions]{Borel graphs generated by commuting functions}
\author{Su Gao, Xiangxi Hu, Jie Zou}
\date{Source: arXiv:2608.15422v2 (CC BY 4.0)}
\begin{document}
\maketitle

\noindent\emph{Source and licence.} Borel graphs generated by commuting functions. Version arXiv:2608.15422v2 (CC BY 4.0), distributed under the Creative Commons Attribution 4.0 International licence, as recorded in the arXiv licence metadata for this exact version and shown on the abstract page https://arxiv.org/abs/2608.15422v2. Full rights notice: licenses/arxiv-2608.15422-CC-BY-4.0.txt.\par\smallskip

\noindent\textit{Editorial scope.} Counted unit: the Lemma whose source label is \texttt{None} in the article (source0.tex lines 1276--1284, source label \texttt{None}), delivered together with the article's own complete original proof of it (source0.tex lines 1286--1383, one \texttt{proof} environment of 98 lines). No mathematics is written, repaired or re-derived: statement and proof are the article's own text line for line. Required context retained uncounted: none. Every definition, display and label of those lines is retained unchanged. Omitted from this package and not redistributed: 1--1275, 1285--1285, 1384--3448. Nothing inside a retained excerpt is omitted or altered: every retained body is delivered exactly as the article's lines are written, so no omission receipt of any kind accompanies this package. Citations: the retained excerpts carry no citation command at all, so no bibliography entry of the article is transcribed and no reference is redistributed. Reference adaptation: none was needed and none was made; every reference inside the delivered unit resolves inside it, so no unresolved reference or citation is produced. Printed slips kept literal: no printed word, symbol, hypothesis or number of the counted statement or of its proof was altered. Layout and class: the article's own class is \texttt{amsart} with packages including graphicx, float, tikz, etoolbox, geometry, amsmath, amssymb, amsthm, mathtools, enumitem, hyperref; the wrapper is the lane's pinned \texttt{amsart} editorial wrapper at 10pt on A4, which restates the theorem-like environments the retained text needs and adds the article's own macro definitions that the retained text needs (none), with extra wrapper packages (bm, bbm, dsfont, xspace, cancel, mathtools). The delivered numbering is the unit's own. Assets: no retained span contains a figure environment, a graphics-inclusion command, a PDF-page inclusion, a table or a TikZ picture; no image file is copied into this package, the package declares no asset dependency and no image file is redistributed with it. Licence: version arXiv:2608.15422v2 of this article is distributed under the Creative Commons Attribution 4.0 International licence, as recorded in the arXiv licence metadata for this exact version and shown on the abstract page https://arxiv.org/abs/2608.15422v2. A full rights notice is delivered at licenses/arxiv-2608.15422-CC-BY-4.0.txt. Characters: every retained excerpt is pure ASCII, so the delivered engine, XeLaTeX with its default Latin Modern text font, reports no missing character.
\begin{lemma}
     \itshape
    % Let $X$ be a standard Borel space. For $n\in \mathbb{N}$, suppose that $f_1,\cdots ,f_n:X\to X$ are countable-to-one commuting Borel functions. 
     Let $r\in\mathbb N^+$. Define a graph $G^r=(F(X),R)$ by
$$
xRy \quad\Longleftrightarrow\quad x\neq y\text{ and }\rho(x,y)\le r.
$$
Then there is a countable proper Borel coloring for $G^r$.
\end{lemma}

\begin{proof} Let $H_r\subseteq F(X)$ be a Borel $r$-forward-independent hitting set for $G_{f_1,\ldots,f_n}\upharpoonright F(X)$.

First consider arbitrary $\alpha_1,\ldots, \alpha_n\in\mathbb N$. Since each $f_i$ is
countable-to-one, the composition
$f_1^{\alpha_1}\cdots f_n^{\alpha_n}$
is countable-to-one. By the Luzin--Novikov theorem, there is a sequence of Borel functions
$\{P_m^{\alpha_1,\ldots, \alpha_n}\}_{m\in\mathbb N}$
such that, for every $z\in F(X)$,
$$
\left(f_1^{\alpha_1}\cdots f_n^{\alpha_n}\right)^{-1}(z)
=
\bigcup_{m\in \mathbb{N}}
P_m^{\alpha_1,\ldots, \alpha_n}(z).
$$

We now define a Borel coloring
$c:F(X)\to\mathbb N^{n+2}$. 
For $x\in F(X)$, let
$c(x)=(j_1,\ldots,j_n,k,\ell)$
where $(j_1+\cdots+j_n,j_1,\ldots,j_n)$ is the least tuple in the
lexicographic order of $\mathbb{N}^{n+1}$ such that
$f_1^{j_1+r}\cdots f_n^{j_n+r}(x)\in H_r$ (the existence of such $j_1,\ldots,j_n$ follows from the hitting property of
$H_r$ applied to
$f_1^r\cdots f_n^r(x)$),
and $k$ and $\ell$ are respectively the least integers such that
$f_1^{j_1}\cdots f_n^{j_n}(x)
\in
P_k^{r,\ldots,r}(H_r)$
and
$P_\ell^{j_1+r,\ldots,j_n+r}
\bigl(
f_1^{j_1+r}\cdots f_n^{j_n+r}(x)
\bigr)
=
x$.



%We first prove the following claim.

\begin{claim}
\label{clm:coloring-separation}
For every $k\in\mathbb N$, if $x\ne y$ and
$x,y\in P_k^{r,\ldots,r}(H_r)$,
then
$\rho(x,y)>r$.
\end{claim}

\noindent \emph{Proof of the claim}.
Toward a contradiction, assume that $\rho(x,y)\le r$. Then there exist
$a_1,\ldots,a_n$, $b_1,\ldots,b_n\in\mathbb N$
with
$a_1+\cdots+a_n+b_1+\cdots+b_n\le r$
such that
$f_1^{a_1}\cdots f_n^{a_n}(x)
=
f_1^{b_1}\cdots f_n^{b_n}(y)$.

Set
$x':=f_1^r\cdots f_n^r(x)$, 
$y':=f_1^r\cdots f_n^r(y)$.
Since $x\neq y$ and $x,y\in P_k^{r,\ldots,r}(H_r)$, we have $x'\neq y'$ and
$x',y'\in H_r$. For each $1\le i\le n$, define
$t_i:=\max\{a_i-b_i,0\}$ 
and
$t'_i:=\max\{b_i-a_i,0\}$.
Then
$t_i t'_i=0$
and
$t_i,t'_i\in [0,r]\cap \mathbb{N}$
for every $1\le i\le n$. A direct computation gives
$f_1^{t_1}\cdots f_n^{t_n}(x')
=
f_1^{t'_1}\cdots f_n^{t'_n}(y')$. However, this contradicts the $r$-forward-independence of $H_r$, since
$x'\ne y'$ are both in $H_r$.
\hfill $\mbox{\qed}_{\mbox{\scriptsize Claim}}$
\medskip


We now verify that $c$ is a proper coloring. For this, take distinct $x,y\in F(X)$ such that
$c(x)=c(y)=(j_1,\ldots,j_n,k,\ell)$.
By the definition of $c$, we have
$f_1^{j_1}\cdots f_n^{j_n}(x),f_1^{j_1}\cdots f_n^{j_n}(y)\in P_k^{r,\ldots,r}(H_r)$.
Moreover, $f_1^{j_1}\cdots f_n^{j_n}(x)\ne f_1^{j_1}\cdots f_n^{j_n}(y)$. Indeed, if $f_1^{j_1}\cdots f_n^{j_n}(x)=f_1^{j_1}\cdots f_n^{j_n}(y)$, then applying
$f_1^r\cdots f_n^r$ would give
$$
f_1^{j_1+r}\cdots f_n^{j_n+r}(x)
=
f_1^{j_1+r}\cdots f_n^{j_n+r}(y),
$$
and the definition of $\ell$ gives $x=y$,
a contradiction. By Claim~\ref{clm:coloring-separation},
$\rho(f_1^{j_1}\cdots f_n^{j_n}(x),f_1^{j_1}\cdots f_n^{j_n}(y))>r$.
Therefore,
$\rho(x,y)>r$.

Thus, $c$ is a proper Borel coloring of $G^r$ with countably many colors. This completes the proof.
\end{proof}

\end{document}
